\section{Conclusion}
\label{sec:conclusion}

The completed F25 reconstruction and extension provide classification evidence
for ten prescribed bridge-damage configurations under operational variation.
In the separately reoptimized F25-X single-channel table, PAA with multi-rate
pooling on carbody pitch rate achieved $77.63\%$ mean test accuracy with a
$2.24$ percentage-point standard deviation over 20 initialization seeds. In the
exploratory frozen-hyperparameter pair table, PAA-CNN using carbody pitch rate
and carbody vertical acceleration achieved a mean of $75.68\%$ with a $4.59$
percentage-point standard deviation. These are descriptive summaries of
different analysis tiers; the latter is not a selected pair for the
four-block regression campaign. Comparisons remain conditional on one
partition, one profile, the specified search budget, and the modeled
observation channels.

F25 is a publication-faithful reconstruction rather than an exact replication.
Its support-aligned mesh, local uniform-$EI$ crack surrogate, software
framework, and explicitly assigned seeds impose documented deviations.
Passage-level test performance measures recognition of the ten fixed
configurations under new operational draws, not generalization to unseen
structural states. The class labels combine three damage axes and do not form
a single ordinal severity scale.

The prospective four-block framework addresses the separate question of
continuous support-stiffness-loss estimation on held-out structural states.
It specifies F40-S, F40-M, L99-S, and L99-M; semantic state grouping and
common-random-number pairing where catalogues overlap; a sealed outer test;
grouped-development out-of-fold adjudication; a fixed channel screen;
independent selected-pair optimization in each block; and disjoint report-only
initialization seeds. F25 results do not complete those analyses.

\pending{Summarize authenticated four-block outcomes after execution: RAW/PAA
factorial winners, selected pair, block-local scour performance and stability,
multi-pier localization, bearing leakage, and paired S--M comparisons.}

Both experiments remain simulation studies. Their modeled response channels
are not validated transducer packages, the wheelset proxies are excluded from
learning, and the stiffness-loss labels are not scour depth. Field deployment
requires measured crossings, a hardware-specific observation model, alternative
track profiles, and explicit evaluation of the simulation-to-field gap.
